\section{Methods}

\subsection{Continual Learning}
\wordcount{250}
\begin{itemize}
    \item Two incremental-learning scenarios are used: Domain-IL, where later tasks share the same output units (e.g.\ relearning road signs in a new country), and Class-IL, where each task adds disjoint output units (e.g.\ learning new animals abroad while still knowing the old ones).
    \item Both scenarios use the same 2-task, 5-class-per-task Split-MNIST protocol, at $14\times14$ resolution, so that the output layer's structure is the only difference between them.
    \item Replay acts as the positive control throughout, establishing that the forgetting problem is solvable in principle under this setup.
\end{itemize}

\subsection{Energy-Based Models (and Backpropagation)}
\wordcount{400}
\begin{itemize}
    \item Predictive coding networks maintain explicit error nodes alongside value nodes, so that error signals are exchanged locally between adjacent nodes rather than backpropagated as a single global signal.
    \item Under prospective configuration, inputs and outputs are clamped and hidden activity is allowed to relax to the state that minimises total prediction error before any weight update is applied.
    \item Backpropagation is run on an identical architecture and objective and serves as the negative control, differing from predictive coding only in how credit is assigned.
\end{itemize}

\subsection{Controls}
\wordcount{250}
\begin{itemize}
    \item Hidden width, architecture, loss, target coding, and learning rate are held identical across rules wherever a fair comparison depends on it; see \cref{sec:results} for which of these are settled and which remain open.
    \item No number from work predating the current experiment series is treated as evidence; every citable claim in this report is re-derived under the present protocol.
    \item Backpropagation and replay are the two controls bracketing the comparison; predictive coding is the rule under test.
\end{itemize}

\subsection{Metrics}

\subsubsection{Endpoint}
\wordcount{150}
\begin{itemize}
    \item The simplest forgetting readout: task-1 accuracy at the end of task-2 training.
    \item It is confounded by the training budget, since every method trends toward the same floor given enough updates, so endpoint alone cannot separate rules reliably.
\end{itemize}

\subsubsection{S\&B}
\wordcount{150}
\begin{itemize}
    \item Song and Bogacz's own combined test-error measure, reported here for comparability with their published claims.
    \item It conflates learning speed with forgetting and cannot separate the two on its own.
\end{itemize}

\subsubsection{Crossover}
\wordcount{200}
\begin{itemize}
    \item The task-1 accuracy at the point where its curve is overtaken by task-2's, measured after the switch.
    \item Invariant to uniform time rescaling, and, empirically, the most stable of the candidate metrics under changes to the stopping rule.
\end{itemize}

\subsubsection{Curve over Diagonal Area}
\wordcount{150}
\begin{itemize}
    \item The integral of the task-1-vs-task-2 trajectory above the line where the two are equal.
    \item A trade-off efficiency measure: how much task-1 accuracy is preserved per unit of task-2 accuracy gained.
\end{itemize}

\subsection{Interpretability}

\subsubsection{Path Efficiency}
\wordcount{200}
\begin{itemize}
    \item Synaptic path length, the accumulated $\sum|\Delta w|$ over training, is compared against the net displacement $|w(T)-w(0)|$ to give a per-weight inefficiency ratio.
    \item A ratio of 1.0 is a straight-line update; a larger ratio means the weight wandered.
\end{itemize}

\subsubsection{Target Alignment}
\wordcount{200}
\begin{itemize}
    \item Song and Bogacz's own proposed mechanism, measured in output space: the cosine similarity between the direction from the current output toward the target, and the direction an update actually moves the output.
    \item Their claim is that predictive coding is more aligned than backpropagation, and that this alignment is what protects task 1.
\end{itemize}

\subsubsection{NCM / Probing}
\wordcount{250}
\begin{itemize}
    \item Nearest-class-mean classification applied directly to hidden-layer features, discarding the output layer entirely.
    \item Separates a genuine representation failure (both the ordinary readout and the probe are low) from a calibration or readout failure (the ordinary readout is low but the probe is high).
\end{itemize}
